\chapter{Установка и первый запуск}

\section{Что именно устанавливать}

Слово \enquote{Docker} сегодня обозначает несколько продуктов. Для обучения
достаточно двух вариантов.

\begin{description}[style=nextline]
  \item[Docker Engine] Демон, CLI и плагины для Linux. Ставится пакетным
        менеджером или официальным скриптом, работает без графической
        оболочки. Это стандартный выбор для серверов и CI.
  \item[Docker Desktop] Приложение для macOS, Windows и Linux с графическим
        интерфейсом, Kubernetes, расширениями и удобным управлением
        ресурсами. Внутри macOS и Windows содержит виртуальную машину
        с Linux, потому что контейнеры используют ядро Linux.
\end{description}

Есть и совместимые альтернативы: Podman, Colima, Rancher Desktop.
Команды почти совпадают, но в этой книге предполагается Docker Engine
или Docker Desktop с CLI \cmd{docker}.

\begin{notebox}
Проверить, что уже установлено, можно командой \cmd{docker version}.
Если команда не найдена~--- Docker нет в системе. Если найдена,
но обращается к серверу с ошибкой~--- установлен только клиент
или не запущен демон.
\end{notebox}

\section{Установка Docker Engine в Linux}

\subsection{Быстрый путь: пакет из репозитория дистрибутива}

В Debian и Ubuntu есть пакет \cmd{docker.io} с демоном и CLI.
Он подходит для первых шагов, но может отставать от актуальной версии
и не включать плагин Compose.

\begin{code}[language=bash]
$ sudo apt-get update
$ sudo apt-get install -y docker.io
$ sudo systemctl enable --now docker
\end{code}

В Fedora и RHEL аналогичную роль играет \cmd{podman-docker} или пакет
\cmd{moby-engine}; для полноценной работы с Compose лучше официальный
репозиторий Docker, описанный ниже.

\begin{tipbox}
Для учебной машины быстрый путь вполне допустим. Для production
используйте официальный репозиторий: вы получите свежие исправления
безопасности и единые версии на всех серверах.
\end{tipbox}

\subsection{Рекомендуемый путь: официальный репозиторий}

Для Ubuntu и Debian установка состоит из четырёх шагов: подключить ключ,
добавить репозиторий, обновить индексы и установить пакеты.

\begin{code}[language=bash]
$ sudo apt-get update
$ sudo apt-get install -y ca-certificates curl
$ sudo install -m 0755 -d /etc/apt/keyrings
$ sudo curl -fsSL https://download.docker.com/linux/ubuntu/gpg \
    -o /etc/apt/keyrings/docker.asc
$ sudo chmod a+r /etc/apt/keyrings/docker.asc
$ echo "deb [arch=$(dpkg --print-architecture) \
    signed-by=/etc/apt/keyrings/docker.asc] \
    https://download.docker.com/linux/ubuntu \
    $(. /etc/os-release && echo $VERSION_CODENAME) stable" \
  | sudo tee /etc/apt/sources.list.d/docker.list > /dev/null
$ sudo apt-get update
$ sudo apt-get install -y docker-ce docker-ce-cli containerd.io \
    docker-buildx-plugin docker-compose-plugin
\end{code}

Для Debian замените \cmd{ubuntu} на \cmd{debian} в URL репозитория,
для других дистрибутивов используйте соответствующие команды
\cmd{dnf} или \cmd{zypper}. Пакеты означают:

\begin{itemize}
  \item \cmd{docker-ce}~--- демон;
  \item \cmd{docker-ce-cli}~--- клиент \cmd{docker};
  \item \cmd{containerd.io}~--- среда выполнения;
  \item \cmd{docker-buildx-plugin}~--- расширенная сборка (BuildKit);
  \item \cmd{docker-compose-plugin}~--- команда \cmd{docker compose}.
\end{itemize}

\begin{warnbox}
Скрипт \cmd{curl -fsSL https://get.docker.com | sh} удобен для быстрых
экспериментов, но не подходит для управляемых серверов: он не фиксирует
версию, не проверяет пакеты и выполняет удалённый код от имени root.
В production устанавливайте пакеты явно.
\end{warnbox}

\section{Проверка установки}

\begin{code}[language=bash]
$ docker version
Client: Docker Engine - Community
 Version:           29.1.5
 ...
Server: Docker Engine - Community
 Engine:
  Version:          29.1.5
  Storage Driver:   overlay2
  Cgroup Driver:    systemd
  Cgroup Version:   2
  ...
\end{code}

Обратите внимание на три поля в секции \cmd{Server}:

\begin{itemize}
  \item \cmd{Storage Driver}~--- драйвер слоистой файловой системы.
        На современных Linux это \cmd{overlay2}.
  \item \cmd{Cgroup Version}~--- версия контрольных групп. Вторая версия
        предпочтительна.
  \item \cmd{Server Version}~--- версия демона. Если секции \cmd{Server}
        нет вовсе, клиент не смог связаться с демоном.
\end{itemize}

Более подробную информацию, включая путь к данным, список плагинов
и настройки ядра, выводит \cmd{docker info}.

\section{Права доступа к демону}

Демон слушает Unix-сокет \file{/var/run/docker.sock}, принадлежащий
пользователю \cmd{root} и группе \cmd{docker}. Поэтому обычный
пользователь получает \enquote{permission denied}. Есть два решения.

\subsection{Добавить пользователя в группу docker}

\begin{code}[language=bash]
$ sudo usermod -aG docker $USER
$ newgrp docker    # применить группу в текущей сессии
\end{code}

После повторного входа в систему проверьте:

\begin{code}[language=bash]
$ docker run --rm hello-world
\end{code}

\begin{warnbox}
Членство в группе \cmd{docker} фактически даёт права root: через демон
можно запустить привилегированный контейнер и смонтировать каталоги
хоста. На общих машинах это неприемлемо. Предпочтительный вариант
для рабочих станций~--- rootless-режим.
\end{warnbox}

\subsection{Rootless-режим}

Rootless-режим запускает демон и контейнеры от имени обычного
пользователя, используя user namespaces. Установка выполняется
отдельным пакетом:

\begin{code}[language=bash]
$ sudo apt-get install -y docker-ce-rootless-extras
$ dockerd-rootless-setuptool.sh install
$ systemctl --user enable --now docker
$ export DOCKER_HOST=unix:///run/user/$(id -u)/docker.sock
\end{code}

Ограничения rootless: нельзя слушать порты ниже 1024 без
дополнительных настроек, некоторые драйверы и функции (например,
некоторые сетевые режимы) недоступны. Зато компрометация контейнера
не даёт root на хосте. Подробнее~--- в разделе~12.3.

\section{Docker Desktop}

В macOS и Windows установите Docker Desktop с официального сайта
и запустите его. Проверьте, что в трее появился значок Docker, а команда
\cmd{docker version} показывает и клиент, и сервер.

Полезные настройки:

\begin{itemize}
  \item \textbf{Resources}~--- сколько процессоров, памяти и диска отдать
        виртуальной машине Docker. Для комфортной работы обычно достаточно
        4~CPU и 8~ГБ памяти.
  \item \textbf{File sharing}~--- какие каталоги хоста доступны для привязок.
        На macOS привязка каталогов вне разрешённых путей может работать
        медленно или не работать вовсе.
  \item \textbf{General $\rightarrow$ Start Docker Desktop when you sign in}~---
        автозапуск, чтобы не забывать про демон.
\end{itemize}

\begin{notebox}
В Windows Docker Desktop использует WSL~2. Держите дистрибутив WSL
на той же файловой системе, где лежит проект: привязка каталогов
с диска \cmd{C:} в WSL заметно медленнее, чем работа внутри дистрибутива WSL.
\end{notebox}

\section{Первый контейнер}

Убедимся, что всё работает, на классическом примере.

\begin{code}[language=bash]
$ docker run --rm hello-world

Unable to find image 'hello-world:latest' locally
latest: Pulling from library/hello-world
...
Hello from Docker!
This message shows that your installation appears to be working correctly.
...
\end{code}

Что произошло по шагам:

\begin{enumerate}
  \item клиент отправил запрос демону;
  \item демон не нашёл образ локально и скачал его из Docker Hub;
  \item из образа был создан контейнер, выполнивший программу \cmd{/hello};
  \item программа напечатала текст и завершилась;
  \item флаг \flag{--rm} удалил контейнер сразу после выхода.
\end{enumerate}

Теперь запустим что-то полезное~--- веб-сервер nginx.

\begin{code}[language=bash]
$ docker run -d --name web -p 8080:80 nginx:alpine
$ docker ps
CONTAINER ID   IMAGE          COMMAND                  CREATED         STATUS         PORTS                  NAMES
a1b2c3d4e5f6   nginx:alpine   "/docker-entrypoint.…"   5 seconds ago   Up 5 seconds   0.0.0.0:8080->80/tcp   web
$ curl -I http://localhost:8080
HTTP/1.1 200 OK
Server: nginx/1.27.3
...
\end{code}

Разберём флаги:

\begin{itemize}
  \item \flag{-d}~--- запустить в фоне (detached);
  \item \flag{--name web}~--- дать контейнеру имя вместо случайного;
  \item \flag{-p 8080:80}~--- опубликовать порт 80 контейнера
        на порту 8080 хоста;
  \item \cmd{nginx:alpine}~--- имя образа и тег.
\end{itemize}

Остановим и удалим контейнер:

\begin{code}[language=bash]
$ docker stop web
$ docker rm web
\end{code}

\begin{tipbox}
Интерактивный эксперимент удобно проводить командой
\cmd{docker run --rm -it alpine sh}: флаг \flag{-i} оставляет stdin
открытым, \flag{-t} выделяет псевдотерминал, \flag{--rm} убирает
контейнер после выхода из оболочки.
\end{tipbox}

\section{Настройка демона}

Основные параметры демона задаются в файле \file{/etc/docker/daemon.json}.
После изменения файла демон нужно перезапустить.

\begin{code}[language=bash]
{
  "log-driver": "json-file",
  "log-opts": {
    "max-size": "10m",
    "max-file": "3"
  },
  "live-restore": true
}
\end{code}

\begin{itemize}
  \item \cmd{log-driver} и \cmd{log-opts} ограничивают размер журналов
        контейнеров. Без ограничения файлы журналов могут занять весь диск.
  \item \cmd{live-restore} позволяет контейнерам пережить перезапуск
        демона (не работает в режиме Swarm).
\end{itemize}

\begin{warnbox}
Не редактируйте \file{daemon.json} вслепую: неверный JSON не даст демону
запуститься. Проверяйте синтаксис, например
\cmd{python3 -m json.tool /etc/docker/daemon.json}, и держите копию
рабочей конфигурации.
\end{warnbox}

\section{Если что-то пошло не так}

\begin{center}
\small
\begin{tabularx}{\textwidth}{@{}XX@{}}
\toprule
\textbf{Симптом} & \textbf{Что проверить} \\
\midrule
\cmd{permission denied ... docker.sock} & пользователь не в группе
\cmd{docker} или не включён rootless-режим \\
\cmd{Cannot connect to the Docker daemon} & демон не запущен:
\cmd{systemctl status docker}; проверьте \envvar{DOCKER\_HOST} \\
\cmd{port is already allocated} & порт занят другим процессом:
\cmd{ss -ltnp} или смените порт хоста \\
\cmd{no space left on device} & \cmd{docker system df}; очистите
неиспользуемые образы и тома \\
\cmd{pull access denied} или превышен лимит & выполните
\cmd{docker login} и проверьте имя образа \\
ошибки \cmd{iptables} при старте & конфликт с firewall или nftables;
смотрите \cmd{journalctl -u docker} \\
\bottomrule
\end{tabularx}
\end{center}

\begin{summarybox}
\begin{itemize}
  \item На Linux ставьте Docker Engine из официального репозитория;
        на macOS и Windows удобнее Docker Desktop.
  \item После установки проверьте \cmd{docker version} и \cmd{docker info}.
  \item Доступ к сокету~--- это фактически root; используйте группу
        \cmd{docker} осознанно или rootless-режим.
  \item Первый контейнер запускается командой \cmd{docker run},
        а \flag{-d}, \flag{-p} и \flag{--name}~--- три флага, которые
        понадобятся чаще всего.
\end{itemize}
\end{summarybox}

\section*{Вопросы для самопроверки}
\begin{selfcheck}
  \item Чем Docker Engine отличается от Docker Desktop?
  \item Почему обычный пользователь не может подключиться к сокету Docker?
  \item Какие три поля вывода \cmd{docker version} стоит проверить сразу
        после установки?
  \item Что делает флаг \flag{--rm} и чем он отличается от \cmd{docker rm}?
  \item Зачем ограничивать размер журналов контейнеров?
\end{selfcheck}
